% Copyright (c) 2026 Geovana Neves. Licensed under CC BY 4.0 (https://creativecommons.org/licenses/by/4.0/). See LICENSE-AND-AUTHORSHIP.md.
%
% 01-fts-overview/main.tex
%
% Guide 1 of the pyflightstream guides: the overview read before guides 2 to 9. What
% FlightStream is and how a surface vorticity panel solver works, the
% simulation workflow from geometry to post-processing, how pyflightstream
% carries that workflow in a workspace, and a map of the other seven guides.
% Built by ../build-all.ps1 (or build-all.sh) with this folder and ../shared/
% as include directories. Its own figures are in tikz/, named g01-*.
%
% No em dashes and no en dashes. No underscores in frame or section titles
% outside \code: beamer writes titles to .nav and .toc.

\documentclass[aspectratio=169,11pt]{beamer}

\input{preamble-slides.tex}
\input{hooks-slides.tex}
\ftsinput{boxes}
\ftsinput{hooks-contract}

\usepackage[nohyphen]{underscore}

\input{info.tex}
\guidepdftitle{1}{Overview}

% Table columns, ragged right so a narrow cell never hyphenates or stretches:
% P{width} for a fixed column, R for the one that takes the rest.
\newcolumntype{P}[1]{>{\RaggedRight\arraybackslash}p{#1}}
\newcolumntype{R}{>{\RaggedRight\arraybackslash}X}

\begin{document}

\guidetitleframe{1}{FlightStream and pyfts, an overview}{Overview}
  {What the solver does, how a simulation is built, and where each guide
   takes over}
  {Sources: the FlightStream user manual, the panel-method literature, and
   the package pages ``The workspace and the workflow'' and ``Storage and
   sync'' [1, 2, 6, 7, 8, 10].}

\begin{frame}{Who reads what}
\begin{columns}[T]
\begin{column}{0.48\textwidth}
\begin{block}{You are new to FlightStream}
\begin{itemize}\small
  \item Part 1: what the tool is, and what it will not do.
  \item Part 2: how a surface vorticity panel solver works, steady,
        unsteady and rotating.
  \item Part 3: the simulation workflow, stage by stage.
\end{itemize}
\end{block}
\end{column}
\begin{column}{0.48\textwidth}
\begin{block}{You know FlightStream, and meet pyfs}
\begin{itemize}\small
  \item Part 4: how pyflightstream carries each stage: the matrix, the
        references, the setups, the pproc, run and post, the four
        workflows, storage and sync.
  \item Part 5: the other seven guides, and when to read each.
\end{itemize}
\end{block}
\end{column}
\end{columns}
\vspace{2mm}
\begin{takeaway}This deck orients; it defines nothing. Every key, command and
refusal it names is defined in the guide it points to.\end{takeaway}
\end{frame}

\begin{frame}{Contents}
\begin{multicols}{2}
\tableofcontents
\end{multicols}
\end{frame}

% Copyright (c) 2026 Geovana Neves. Licensed under CC BY 4.0 (https://creativecommons.org/licenses/by/4.0/). See LICENSE-AND-AUTHORSHIP.md.
%
% 01-fts-overview/text/01-flightstream.tex
%
% Part 1: what FlightStream is, where it sits on the fidelity ladder, what it
% will not do, why the build is part of every answer, and the published
% record of the method. Manual pages are those of build 26121 [6].

\section{FlightStream}

\begin{frame}{Where a panel method sits}
\begin{columns}[T]
\begin{column}{0.46\textwidth}
\begin{itemize}\small
  \item Aerodynamic tools form a \textbf{fidelity ladder}: each level resolves
        more physics and costs more per case.
  \item A panel method sits in the middle: the \textbf{whole configuration} in
        three dimensions, in seconds to minutes per point on a workstation.
  \item That is the regime of \textbf{interference and installation}: where a
        propeller sits, what a fuselage does to a wing, a pitching moment slope.
\end{itemize}
\end{column}
\begin{column}{0.52\textwidth}
\slidefig[0.52]{g01-fidelity-ladder}
\end{column}
\end{columns}
\begin{takeaway}A panel method buys full-configuration aerodynamics at a cost
that still allows design iteration.\end{takeaway}
\end{frame}

\begin{frame}{What FlightStream is}
\begin{columns}[T]
\begin{column}{0.54\textwidth}
\begin{itemize}\small
  \item A \textbf{surface vorticity} panel solver: the potential-flow problem,
        with vorticity carried on the surface mesh [7, 8]. It tolerates
        unstructured, engineering-quality meshes and reads the common exchange
        formats [6, pp.~72-75].
  \item The core is \textbf{inviscid}. The rest is correction models on top of
        it: a boundary layer, with optional viscous coupling, for friction drag
        and decambering [6, pp.~204-211]; subsonic compressibility, chosen at
        initialisation [6, pp.~218-221], [12].
  \item Steady, unsteady and rotating solutions, driven from a graphical
        interface or from a \textbf{script} of commands.
\end{itemize}
\end{column}
\begin{column}{0.42\textwidth}
\begin{alertblock}{Out of scope by construction}
\begin{itemize}\small
  \item Massively separated flow, deep stall, buffet.
  \item Transonic shocks.
  \item Acoustics.
\end{itemize}
\end{alertblock}
\begin{alertblock}{And it does not warn you}
\small A case that leaves the validity envelope still converges and still
prints coefficients.
\end{alertblock}
\end{column}
\end{columns}
\end{frame}

\begin{frame}{The build number is part of the answer}
\begin{itemize}\small
  \item FlightStream is released as \textbf{builds}; this series names build
        \fsbuild{} in its examples. The scripting language changes between
        builds of one version: a command that runs on one build may be absent,
        renamed or spelled differently on another.
  \item The manual ships inside each build, and \textbf{editions are not page
        aligned}. A page number is quoted with its edition, or not at all.
  \item A coefficient without the build that produced it is \textbf{not
        reproducible}. Record the build next to every result.
\end{itemize}
\begin{block}{How pyflightstream holds this}
\small Every command it writes comes from a \textbf{command database per
build}, and every row of a study names its build in the column
\code{FS\_BUILD}. \pyfs{plan} checks the whole study against that build's
database before any seat is spent (Part 4, Guide 2).
\end{block}
\end{frame}

\begin{frame}{The tool's own view of the workflow}
\begin{columns}[T]
\begin{column}{0.50\textwidth}
\begin{itemize}\small
  \item The manual closes with three checklists: the \textbf{general}
        workflow [6, p.~411], the workflow from \textbf{CAD cross sections}
        (CCS) [6, p.~412], and the \textbf{rotor and turbine} workflow
        [6, p.~413].
  \item The general one has five steps: geometry preparation, half-model
        setup, surface detection, boundary conditions, simulation setup.
  \item Boundary conditions come \textbf{before} anything a user would call
        setup. Part 3 explains why: initialisation locks them.
\end{itemize}
\end{column}
\begin{column}{0.46\textwidth}
\begin{block}{A half model is expected}
\small The manual recommends the half model under a mirror plane, and free
edges along that plane are \textbf{expected}, not a defect [6, p.~411].
\end{block}
\begin{block}{A reference case ships with it}
\small A rotor in hover and cruise [6, p.~401]: the fastest check that an
installation behaves.
\end{block}
\end{column}
\end{columns}
\end{frame}

\begin{frame}{The method, and where it has been measured}
\begin{columns}[T]
\begin{column}{0.48\textwidth}
\begin{block}{What it solves}
\begin{itemize}\small
  \item The integrated-circulation formulation, in full [7].
  \item The surface vorticity panel method in the archival literature: the
        reference to cite when someone asks what FlightStream solves [8].
  \item The classical panel-method background: [9, 10].
\end{itemize}
\end{block}
\end{column}
\begin{column}{0.48\textwidth}
\begin{block}{Where it has been compared with experiment}
\begin{itemize}\small
  \item Stability and control of a vehicle against wind-tunnel data:
        derivatives agree acceptably, control effectiveness is overpredicted,
        propellers off [13].
  \item A high-lift wing with distributed propulsion, against experiment: the
        closest published case to a powered airframe [14].
\end{itemize}
\end{block}
\end{column}
\end{columns}
\vspace{1mm}
\begin{takeaway}Convergence against a theoretical anchor is
\emph{verification}; agreement with measured data is \emph{validation}. Say
which one a result shows.\end{takeaway}
\end{frame}

% Copyright (c) 2026 Geovana Neves. Licensed under CC BY 4.0 (https://creativecommons.org/licenses/by/4.0/). See LICENSE-AND-AUTHORSHIP.md.
%
% 01-fts-overview/text/02-theory.tex
%
% Part 2: the theory a user needs to justify the boundary-condition and setup
% choices, steady, unsteady and rotating, and nothing more. Manual pages are
% those of build 26121 [6]; the physics is in [7-12, 15-17].

\section{How a panel solver works}

\begin{frame}{Potential flow and its boundary conditions}
\begin{columns}[T]
\begin{column}{0.56\textwidth}
\begin{itemize}\small
  \item Assumptions: \textbf{incompressible, irrotational, inviscid}. The
        velocity derives from a potential, \(\mathbf{V} = \nabla\phi\), and
        mass conservation becomes the Laplace equation,
        \(\nabla^2\phi = 0\) [9, 10].
  \item The equation is \textbf{linear}: elementary solutions superpose, and
        the flow about a body is assembled from them.
  \item Two conditions close it: \textbf{flow tangency} on the surface, and
        \textbf{decay} of the disturbance far away.
\end{itemize}
\end{column}
\begin{column}{0.40\textwidth}
\begin{block}{What is absent}
\small No volume grid, so \textbf{no farfield to size} and no blockage to
correct: the disturbance decays because the elementary solutions decay
[6, p.~164].
\end{block}
\begin{block}{What it can predict}
\small Pressure effects: lift, induced drag, moments, interference. Never
friction drag, which lives in a boundary layer it does not model.
\end{block}
\end{column}
\end{columns}
\end{frame}

\begin{frame}{From source and doublet panels to surface vorticity}
\begin{columns}[T]
\begin{column}{0.58\textwidth}
\begin{itemize}\small
  \item A classic panel method gives each surface panel a source or doublet
        strength, enforces tangency at one point per panel, and solves a
        dense linear system [9, 10].
  \item FlightStream carries \textbf{vorticity} on the surface mesh instead
        [7, 8]. Degenerate faces are removed in preprocessing, and
        flow-aligned quadrilaterals are preferred on lifting surfaces
        [6, pp.~75, 156].
  \item Whatever the singularity: \textbf{surface resolution drives the
        induced velocities}, and cost grows steeply with the panel count.
\end{itemize}
\end{column}
\begin{column}{0.38\textwidth}
\slidefig[0.46]{g01-panels-and-wake}
\end{column}
\end{columns}
\begin{takeaway}Different singularity, same physics: enforce tangency, resolve
the circulation, integrate the pressures.\end{takeaway}
\end{frame}

\begin{frame}{The Kutta condition and the wake}
\begin{itemize}\small
  \item Potential flow alone \textbf{does not fix lift}: any circulation
        satisfies the equations. The physical statement that selects the real
        solution is smooth flow leaving a sharp trailing edge [10, 11].
  \item In the solver, \textbf{marking a trailing edge} is that statement:
        wake filaments are shed from marked edges, and an unmarked edge sheds
        nothing at all [6, p.~165].
  \item The wake carries the trailed vorticity downstream. Its strength
        mirrors the spanwise load gradient, and \textbf{induced drag} is the
        price of that vorticity.
\end{itemize}
\begin{alertblock}{The most critical input}
\small An unmarked trailing edge raises no error. It returns a converged
solution with the wrong lift; the manual calls the trailing edges by far the
most critical input [6, p.~164].
\end{alertblock}
\end{frame}

\begin{frame}{How much circulation is shed: the factor \(K\)}
\begin{itemize}\small
  \item The strength of the wake filament leaving a marked edge is
        \[ \Gamma_{s_t} = \tfrac{K}{2}\left(q_u^{2} - q_l^{2}\right)\Delta t, \]
        with \(q_u\) and \(q_l\) the velocities above and below the filament
        [6, p.~165].
  \item \textbf{Standard Kutta}, \(K = 1\): the full net circulation is shed.
        The default, and the right choice at any real lifting trailing edge.
  \item \textbf{Relaxed Kutta}, \(K = 0.5\): half is shed. For an edge that is
        not sharp, typically the aft end of a rounded fuselage, so that the
        body carries an equivalent lifting load rather than none.
  \item Two more types exist: exhaust outflow, and vortex shedding, which has
        no vortex-bursting model and is not recommended beyond 25 to 30 degrees
        of incidence [6, p.~165].
\end{itemize}
\begin{takeaway}Standard against relaxed is not a numerical knob: it states
whether an edge is a real trailing edge.\end{takeaway}
\end{frame}

\begin{frame}{When a trailing edge is not a Kutta edge}
\begin{columns}[T]
\begin{column}{0.54\textwidth}
\begin{itemize}\small
  \item A \textbf{blunt} trailing edge is a separated base: the solver applies
        a constant pressure over a \textbf{base region} [6, p.~164].
  \item Rule of thumb: thicker than \textbf{2 percent of the chord}, keep it
        blunt with a base region, and mark the base perimeter as trailing
        edges [6, pp.~164, 167]. Thinner, collapse it to a sharp edge: no loss
        of accuracy, faster and cleaner convergence.
  \item A body protruding through a wing trailing edge by a length \(L\), on a
        chord \(c\): the carry-over decision turns on \(L > c/2\)
        [6, p.~166].
\end{itemize}
\end{column}
\begin{column}{0.42\textwidth}
\begin{block}{Taken in CAD}
\small Collapsed or blunt is decided \textbf{before the mesh is made}, and it
governs every stage after it. Part 3 places it in the geometry stage.
\end{block}
\end{column}
\end{columns}
\end{frame}

\begin{frame}{What the method captures, and what it does not}
\begin{columns}[T]
\begin{column}{0.50\textwidth}
\begin{itemize}\small
  \item \textbf{Inviscid core}: lift slope, induced drag, spanwise loading,
        thickness effects, interference between components.
  \item \textbf{Correction models}: friction drag by the boundary layer,
        decoupled by default, coupled through transpiration for nonlinear
        decambering [6, pp.~204-211]; subcritical Mach effects by
        Prandtl-Glauert [6, pp.~218-221], [12].
  \item \textbf{Not captured} under any setting: deep stall, buffet, shock
        stall.
\end{itemize}
\end{column}
\begin{column}{0.46\textwidth}
\slidefig[0.40]{g01-trust-region}
\end{column}
\end{columns}
\begin{takeaway}A polar is trustworthy where the flow is attached, and nothing
in the output says where that ends.\end{takeaway}
\end{frame}

\begin{frame}{Unsteady: the wake remembers}
\begin{itemize}\small
  \item \textbf{Kelvin}: in inviscid flow the circulation of a material
        contour is conserved. When the bound circulation changes, an equal and
        opposite vortex is shed into the wake [10].
  \item The \textbf{starting vortex} is the canonical picture: a wing set
        impulsively in motion leaves behind a vortex that mirrors the
        circulation it gains.
  \item The wake stops being a fixed sheet and becomes part of the solution:
        it carries the history of every change of lift, and induces velocity
        back on the wing.
\end{itemize}
\slidefig[0.34]{g01-starting-vortex}
\end{frame}

\begin{frame}{The time step is a mesh dimension}
\begin{itemize}\small
  \item The solver marches in \textbf{physical time}: each step runs a block of
        iterations under the steady settings, and the wake grows one shed
        segment per step [6, pp.~212-213].
  \item The streamwise length of a shed segment is the convection velocity
        times the step: \textbf{the step sets the wake resolution} as the panel
        size sets the surface resolution. Choose it from a physical time the
        geometry has, for example the chord's convective time \(t_c = c/V\).
  \item The judge is the \textbf{reduced frequency} [15, 10]:
        \[ k = \frac{\omega c}{2 V_\infty}. \]
        At low \(k\) the flow adjusts quasi-statically; as \(k\) grows, wake
        memory and phase lag become the answer.
\end{itemize}
\begin{takeaway}Run unsteady only when the question demands it: gusts,
oscillation, dynamic derivatives, rotors. A slow angle sweep is a sequence of
steady points.\end{takeaway}
\end{frame}

\begin{frame}{Rotation added to the unsteady problem}
\begin{itemize}\small
  \item Each blade section sees the sum of \(\Omega r\) and the free stream;
        the inflow angle \(\varphi = \arctan\big(V_\infty / (\Omega r)\big)\)
        varies along the radius, which is why blades are twisted [16, 17].
  \item Blade geometry and section loads live in the \textbf{rotating} frame;
        the wake, the airframe and the performance numbers in the
        \textbf{inertial} one. Name the frame when you report a number.
  \item Everything unsteady still holds. Rotation adds geometry to the wake,
        not new principles.
\end{itemize}
\slidefig[0.36]{g01-velocity-triangles}
\end{frame}

\begin{frame}{Propeller performance, and the helical wake}
\begin{columns}[T]
\begin{column}{0.48\textwidth}
\begin{align*}
J &= \frac{V_\infty}{nD}, & C_T &= \frac{T}{\rho n^2 D^4}, \\
C_P &= \frac{P}{\rho n^3 D^5}, & \eta &= \frac{J C_T}{C_P}
\end{align*}
{\small \(n\) in revolutions per second, \(D\) the diameter [16, 17]. At
fixed \(J\) the non-dimensional picture is fixed; the rotation rate alone only
rescales the Reynolds and tip Mach numbers.\par}
\vspace{1mm}
\begin{alertblock}{The solver reports forces}
\small It reports forces and moments normalised on the reference values you
set; \(J\), \(C_T\), \(C_P\) and \(\eta\) are built by you, or by the
post-processing (Guide 7).
\end{alertblock}
\end{column}
\begin{column}{0.48\textwidth}
\slidefig[0.28]{g01-helix-pitch}
\begin{itemize}\small
  \item The wake is a \textbf{helix} whose pitch is set by \(J\): a low \(J\)
        packs the turns behind the disc, and each blade meets the wake of the
        others.
  \item The clock is the shaft: resolution in degrees per step, about 10 to 12
        degrees, run length in revolutions, reference velocity the tip speed
        [6, p.~413].
\end{itemize}
\end{column}
\end{columns}
\end{frame}

% Copyright (c) 2026 Geovana Neves. Licensed under CC BY 4.0 (https://creativecommons.org/licenses/by/4.0/). See LICENSE-AND-AUTHORSHIP.md.
%
% 01-fts-overview/text/03-workflow.tex
%
% Part 3: the simulation workflow as ten stages, from the question to the
% engineering analysis, with what each stage decides. Manual pages are those
% of build 26121 [6].

\section{The simulation workflow}

\begin{frame}{Ten stages, in the order decisions constrain each other}
\relax
{\ftstablesize\renewcommand{\arraystretch}{1.12}
\begin{tabularx}{\linewidth}{@{}r@{\hspace{2mm}}P{0.34\linewidth}R@{}}
\toprule
\textbf{\#} & \textbf{Stage} & \textbf{The decision it owns} \\
\midrule
1  & Case study and scope & what question the case answers, and what it does not \\
2  & Matrix definition & which points are run, and why those \\
3  & Water-tight geometry & blunt or sharp trailing edge, and where the components split \\
4  & Surface mesh & where resolution is spent \\
5  & Boundary conditions & which edges shed vorticity, and how much \\
6  & Solver setup & what the solver is allowed to assume \\
7  & Execution and monitoring & whether the answer has converged \\
8  & Output definition and export & what evidence the run leaves behind \\
9  & Verification and validation & whether the answer is right \\
10 & Engineering analyses & what the answer means for the design \\
\bottomrule
\end{tabularx}}
\vspace{1mm}
\begin{takeaway}The stages are not a checklist to tick: a decision taken out of
order has to be taken again.\end{takeaway}
\end{frame}

\begin{frame}{One session in the interface, and the lock in it}
\slidefig[0.34]{g01-workflow-chart}
\begin{itemize}\small
  \item \textbf{Initialisation is a lock.} It freezes the solver against the
        mesh and the boundary conditions as they stand, and fixes the solver
        mode and the symmetry claim [6, pp.~218-221].
  \item A runtime setting (angle, velocity, iterations) can change after it.
        A mesh or boundary change needs the initialisation removed and taken
        again. That is why boundary conditions precede setup [6, p.~411].
\end{itemize}
\end{frame}

\begin{frame}{Stages 3 and 4: geometry and surface mesh}
\begin{columns}[T]
\begin{column}{0.48\textwidth}
\begin{block}{Geometry}
\begin{itemize}\small
  \item \textbf{Water-tight} means the body closes, not the domain.
  \item Collapsed or blunt trailing edge: decided in CAD (Part 2).
  \item Components are split by the \textbf{boundary condition they will
        receive}, and by the loads you will want per component, never by CAD
        convenience.
\end{itemize}
\end{block}
\end{column}
\begin{column}{0.48\textwidth}
\begin{block}{Surface mesh}
\begin{itemize}\small
  \item Four sources: the tool's cross-section meshers, transfinite
        interpolation, the wrapper, or an imported mesh [6, pp.~72-75].
        Only the first two regenerate in seconds.
  \item Faces go where the gradients are: leading edge, tip, junction, the
        row the wake leaves from [6, p.~62].
  \item A mesh convergence study is an obligation, and it is affordable.
\end{itemize}
\end{block}
\end{column}
\end{columns}
\begin{takeaway}The mesh carries the boundary conditions. A boundary is a
decision about what you will be able to ask for later.\end{takeaway}
\end{frame}

\begin{frame}{Stage 5: boundary conditions are marking}
\begin{columns}[T]
\begin{column}{0.50\textwidth}
\begin{itemize}\small
  \item \textbf{Farfield}: there is none to set. \textbf{Walls}: every face is
        a slip wall once meshed, switched to no slip by the solver when the
        boundary layer is active [6, p.~164].
  \item What remains is \textbf{marking} where that default is wrong:
        trailing edges, with their type (Part 2); base regions; wake
        termination nodes; and, when the case needs them, actuator discs,
        inlets and outlets, custom free-stream fields [6, pp.~164-197].
\end{itemize}
\end{column}
\begin{column}{0.46\textwidth}
\begin{alertblock}{Auto-detection is a proposal}
\small The solver can detect trailing edges; it proposes, you decide. Inspect
every marked edge before initialisation: a split trailing edge is two edges,
and a missing one is silent.
\end{alertblock}
\end{column}
\end{columns}
\end{frame}

\begin{frame}{Stage 6: what the solver may assume}
\relax
{\ftstablesize\renewcommand{\arraystretch}{1.12}
\begin{tabularx}{\linewidth}{@{}P{0.14\linewidth}R@{}}
\toprule
\textbf{Symmetry} & \textbf{What it asserts} \\
\midrule
None & nothing: the mesh in front of the solver is the whole problem \\
Mirror & a symmetry plane in the flow, not only in the geometry; false as
  soon as the sideslip is not zero \\
Periodic & an onset flow the same at every azimuth, plus a copy count; false
  as soon as the angle of attack is not zero \\
\bottomrule
\end{tabularx}}
\vspace{1mm}
\begin{itemize}\small
  \item \textbf{Coordinate systems first}: loads are reported in one, and a
        moment without its frame and point is not a number.
  \item \textbf{Reference velocity, area and length} normalise every
        coefficient. Set them; never inherit them.
  \item \textbf{Convergence is three controls}: an iteration cap, a residual
        threshold, and how long the residual must stay below it
        [6, pp.~202-203]. Viscous settings change the answer; leaving them off
        is a decision too.
\end{itemize}
\end{frame}

\begin{frame}{Stages 7 and 8: run, watch, and declare the outputs}
\begin{columns}[T]
\begin{column}{0.48\textwidth}
\begin{block}{Execution and monitoring}
\begin{itemize}\small
  \item \textbf{Converged is a state that was held}, not a number that was
        touched. A run stopped at its cap ran out of iterations.
  \item Unsteady asks twice: the residual inside each step, and the history
        of the quantity you want. They are independent.
  \item Periodic: judge the \textbf{per-revolution} statistic, not the
        sample.
\end{itemize}
\end{block}
\end{column}
\begin{column}{0.48\textwidth}
\begin{block}{Output definition and export}
\begin{itemize}\small
  \item \textbf{An output not declared was not computed}: sections, probes
        and plots are set before the run.
  \item Symmetry-load doubling changes what the solver computes. Two runs set
        differently are not the same quantity.
  \item Every coefficient carries a normalisation: record it with the number.
\end{itemize}
\end{block}
\end{column}
\end{columns}
\begin{takeaway}Recovering a variable that was never declared means running
the case again, which for a rotor is the most expensive mistake
available.\end{takeaway}
\end{frame}

\begin{frame}{Stages 9 and 10: is it right, and what does it mean}
\begin{itemize}\small
  \item \textbf{Verification}: the solution against a theoretical anchor and
        against itself. A finite-wing lift slope, a span efficiency, a mesh
        and a time-step study, a sector against the full wheel.
  \item \textbf{Validation}: the solution against measured data, inside the
        envelope where the method has been compared with experiment
        [13, 14].
  \item An error that scales the whole solution divides out of a ratio. An
        efficiency can agree while thrust and torque are both off: check the
        forces, not only their ratio.
  \item \textbf{Engineering analyses} use the answer: derivatives, trim,
        installation deltas, component breakdowns. They inherit every
        assumption of stages 3 to 8.
\end{itemize}
\begin{takeaway}A mid-fidelity number is defensible when its build, its
settings, its convergence and its verification travel with it.\end{takeaway}
\end{frame}

% Copyright (c) 2026 Geovana Neves. Licensed under CC BY 4.0 (https://creativecommons.org/licenses/by/4.0/). See LICENSE-AND-AUTHORSHIP.md.
%
% 01-fts-overview/text/04-pyfs.tex
%
% Part 4: how pyflightstream carries the workflow of Part 3 in a workspace:
% the matrix, the three input files a row cites, the four stages of a run,
% the four workflows, and storage and sync. Every statement here is defined
% in Guides 2 to 9 and in the package pages [1-4]; this part only places it.

\section{pyflightstream on the workflow}

\begin{frame}{What pyflightstream is}
\begin{columns}[T]
\begin{column}{0.52\textwidth}
\begin{itemize}\small
  \item A Python driver for FlightStream, MIT licensed, installed from PyPI
        (\code{pip install pyflightstream}). FlightStream itself, and its
        licence, are separate.
  \item It \textbf{writes} the solver script from a command database per
        build, \textbf{runs} it, \textbf{reads} what the solver exported, and
        \textbf{post-processes} it into tables with their provenance.
  \item A study is a \textbf{workspace}: a folder with a run matrix and the
        input files the matrix cites. No Python is written to run it.
\end{itemize}
\end{column}
\begin{column}{0.44\textwidth}
\slidefig[0.56]{g01-pyfs-layers}
\end{column}
\end{columns}
\begin{takeaway}A row names things, and the package makes the link: a row never
carries a path.\end{takeaway}
\end{frame}

\begin{frame}{The workspace}
\slidefig[0.80]{guide-workspace-tree}
\end{frame}

\begin{frame}{The ten stages, and where each one lives}
\relax
{\ftstablesize\renewcommand{\arraystretch}{1.08}
\begin{tabularx}{\linewidth}{@{}P{0.25\linewidth}RP{0.10\linewidth}@{}}
\toprule
\textbf{Stage} & \textbf{In the workspace} & \textbf{Guide} \\
\midrule
1, 2 Scope and matrix & one row of \code{matriz.fs} per study point or sweep:
  \code{DESCRIPTION}, \code{FLIGHT\_CONDITION}, \code{SWEEP\_VALUES},
  \code{RUN} & 1 \\
3, 4 Geometry and mesh & made outside the package; the file sits in
  \code{inputs/geometries/}, the row names it in \code{GEOMETRY};
  \pyfs{pyfs-matrix inventory} reads its boundaries & 1, 2 \\
5 Boundary conditions & the trailing edge by file or by detection, base
  regions and wake nodes as keys; rotors and actuator discs in the
  reference file & 2, 3 \\
6 Solver setup & \code{inputs/setups/s<id>.toml} (column \code{SET}); lengths,
  moment point and frames in \code{inputs/references/r<id>.toml} (\code{REF});
  \code{SYMMETRY} and \code{WORKFLOW} on the row & 3, 4 \\
7 Execution & \pyfs{plan}, \pyfs{run}, \pyfs{collect}; \code{FS\_BUILD},
  \code{NCPUS}, \code{WALLTIME} on the row & 1, 7 \\
8 Output definition & \code{inputs/pproc/p<id>.toml} (\code{PPROC}): what the
  run exports, declared before it runs & 4, 5 \\
9, 10 Verification and analyses & \pyfs{post} writes \code{post/}: polars,
  series, sections, probes, provenance; FSI couples a structure & 5, 6 \\
\bottomrule
\end{tabularx}}
\end{frame}

\begin{frame}{The row names things, the files define them}
\slidefig[0.60]{guide-inputs-into-run}
\begin{takeaway}\code{REF}, \code{SET} and \code{PPROC} carry a file stem and
\code{GEOMETRY} names the mesh: stated once, in one file, and every row citing
it demonstrably used it.\end{takeaway}
\end{frame}

\begin{frame}{Reference, setup, pproc: three files a row cites}
\begin{columns}[T]
\begin{column}{0.32\textwidth}
\begin{block}{\code{r<id>.toml}}
\small \textbf{What the coefficients divide by, and about which point}:
reference area and lengths, the moment point, aliases for groups of
surfaces, custom frames, and one block per rotor or actuator disc
(Guide 5).
\end{block}
\end{column}
\begin{column}{0.32\textwidth}
\begin{block}{\code{s<id>.toml}}
\small \textbf{How hard to solve and when to stop}: iterations and
convergence, the boundary layer and its coupling, the model and its
initialisation, wake settings, and a raw table for any command without a key
(Guide 6).
\end{block}
\end{column}
\begin{column}{0.32\textwidth}
\begin{block}{\code{p<id>.toml}}
\small \textbf{What the run exports, and how post reads it}: loads, sections,
probes, plots, the log last; the groups and derived columns post reduces
with. The averaging window belongs to the row (Guides 6, 7).
\end{block}
\end{column}
\end{columns}
\vspace{2mm}
\begin{itemize}\small
  \item The flight condition is a cell of the row, not a file: Mach, speed,
        Reynolds number, altitude, a temperature offset, each a constraint on
        one flow state [1].
  \item Any key only some rows need goes in the row's free cell,
        \code{VAR\_NAMES\_VALUES}.
\end{itemize}
\end{frame}

\begin{frame}[fragile]{Four stages of a run}
\slidefig[0.40]{guide-stage-flow}
\begin{lstlisting}[style=ftsbase]
pyfs-matrix plan    matriz.fs --workspace . --fs-version 26.123   # no seat, no queue
pyfs-matrix run     matriz.fs --workspace . --fs-version 26.123   # solves, or submits
pyfs-matrix collect --workspace .                                 # cluster: settled outputs
pyfs-matrix post    matriz.fs --workspace .                       # rebuilds the products
\end{lstlisting}
{\small \pyfs{plan} validates every row against the build's command database
and writes the receipt \pyfs{run} refuses to start without (Guide 2).\par}
\end{frame}

\begin{frame}{Four workflows, one per kind of physics}
\relax
{\ftstablesize\renewcommand{\arraystretch}{1.10}
\begin{tabularx}{\linewidth}{@{}P{0.17\linewidth}P{0.27\linewidth}RP{0.12\linewidth}@{}}
\toprule
\textbf{\code{WORKFLOW}} & \textbf{What is solved} & \textbf{The physics of Part 2} & \textbf{Cost} \\
\midrule
\code{steady} & one converged steady solve per point; a sweep is one job &
  the steady problem: tangency, Kutta, the fixed wake & seconds per point \\
\code{unsteady} & a march in physical time, nothing turning & Kelvin and
  the starting vortex; the step as a wake dimension; a window to average &
  minutes \\
\code{unsteady\_rotor} & the blades turn, the wake is shed in time & the
  helical wake; degrees per step, revolutions, a per-revolution window; the
  reference for any rotor & minutes to hours \\
\code{qsteady\_rotor} & blades held still, the free stream turning about the
  shaft; sector or wheel & an isolated, axisymmetric rotor where the reduced
  frequency is small & seconds \\
\bottomrule
\end{tabularx}}
\vspace{1mm}
\begin{itemize}\small
  \item The row names its workflow; a key another workflow owns is refused by
        name at \pyfs{plan}. Costs are orders of magnitude on one
        workstation [20].
  \item Guide 2 carries the decision chart, the quasi-steady rotor in depth,
        and when its answer stands for the turning blade.
\end{itemize}
\end{frame}

\begin{frame}{Storage and sync}
\begin{columns}[T]
\begin{column}{0.48\textwidth}
\begin{block}{What a run leaves}
\begin{itemize}\small
  \item \code{sims/}: one folder per matrix row, \code{sim\_<POL>}, with the
        scripts and one datapoint folder per point: the saved simulation and
        the solver's exports.
  \item \code{post/}: the products, rebuilt by \pyfs{post} with no seat.
  \item \code{archive/}: whatever a rebuild replaced.
  \item \code{runs.json}: the record of every run, with its state.
\end{itemize}
\end{block}
\end{column}
\begin{column}{0.48\textwidth}
\begin{block}{Two workspaces, one main}
\begin{itemize}\small
  \item A workstation workspace for planning, one point, looking at results;
        a cluster workspace for long marches and many points.
  \item \pyfs{pyfs-matrix sync}, run on the main, brings the cluster's runs
        home at a chosen level; it previews unless told \code{--apply}, and
        never copies a mesh [2].
\end{itemize}
\end{block}
\end{column}
\end{columns}
\begin{takeaway}You never edit anything under \code{sims/} or \code{post/}. If a
number there is wrong, an input is wrong.\end{takeaway}
\end{frame}

% Copyright (c) 2026 Geovana Neves. Licensed under CC BY 4.0 (https://creativecommons.org/licenses/by/4.0/). See LICENSE-AND-AUTHORSHIP.md.
%
% 01-fts-overview/text/05-guides.tex
%
% Part 5: the map of the other eight guides, the order to read them in for
% four kinds of reader, the public geometries to practise on, and how to cite
% the tool.

\section{The guides}

\begin{frame}{Nine guides, and when to read each}
\relax
{\ftstablesize\renewcommand{\arraystretch}{1.08}
\begin{tabularx}{\linewidth}{@{}P{0.04\linewidth}P{0.31\linewidth}RP{0.24\linewidth}@{}}
\toprule
& \textbf{File} \code{pyfts-guide-} & \textbf{What it covers} & \textbf{Read it when} \\
\midrule
1 & \code{01-fts-overview} & the solver, the workflow, the map &
  before anything else \\
2 & \code{02-workspaces} & the tree, the inputs, choosing a workflow,
  running, two workspaces and sync & you set up or run a study \\
3 & \code{03-gui-to-pyfs} & every step of a GUI session, and the
  key that takes it & you know the interface \\
4 & \code{04-cheatsheet} & every command and option, then one page
  per stage & you need a command, not the explanation \\
5 & \code{05-references} & lengths, the moment point, aliases,
  frames, rotors, actuator discs & you write a reference file \\
6 & \code{06-solver-setup} & the setup file, its keys and the
  solver commands behind them & you tune how a point is solved \\
7 & \code{07-pproc-definitions} & what every product, reduction
  and window means & you read or plot a result \\
8 & \code{08-fsi} & the fluid-structure loop, the structural
  input, calibration, limits [22] & a structure deforms under load \\
9 & \code{09-python-environment-offline} & Python, a virtual
  environment and the package with no internet & once per offline machine \\
\bottomrule
\end{tabularx}}
\end{frame}

\begin{frame}{Four ways through them}
\begin{columns}[T]
\begin{column}{0.48\textwidth}
\begin{block}{A first study}
\small 1, then 2 (Parts 1 to 11), then 5 and 6 for the files you cite, then 7
before you plot, with 4 beside you. The package page ``Getting started'' is the shortest path to
a first run [21].
\end{block}
\begin{block}{From the interface to the workspace}
\small 1 (Parts 3 and 4), then 3 side by side with your own session, then 6.
\end{block}
\end{column}
\begin{column}{0.48\textwidth}
\begin{block}{A rotor}
\small 1 (Part 2, rotation), 5 (the rotor block), 2 (choosing a workflow and
the quasi-steady rotor), 7 (per-blade, phase-locked and window products); 8
for a blade that deforms.
\end{block}
\begin{block}{Installing and operating}
\small 9 on each machine without internet, then 2 (the cluster, two
workspaces and sync, space on disk).
\end{block}
\end{column}
\end{columns}
\begin{takeaway}Each guide is written for one version of the package; its
title page says which.\end{takeaway}
\end{frame}

\begin{frame}{Where to practise, and how to cite}
\begin{columns}[T]
\begin{column}{0.48\textwidth}
\begin{block}{Two public geometries}
\begin{itemize}\small
  \item \textbf{SMI}: an airframe geometry, wing and wing-body, published
        for practice [18].
  \item \textbf{TUD-XPROP}: a research propeller of Delft University of
        Technology, requested from its own record and cited by it [19].
\end{itemize}
\end{block}
\begin{block}{Where to take a question}
\small Solver behaviour to the vendor, with the build and the page. The
driver to the repository, as an issue.
\end{block}
\end{column}
\begin{column}{0.48\textwidth}
\begin{block}{Citing the tool}
\small \pkgcitation
\end{block}
\begin{block}{The Python library}
\small For scripts written in Python rather than a matrix: the library
guide, \code{pyflightstream\_user\_guide.tex} in the repository's
\code{guide/} folder.
\end{block}
\end{column}
\end{columns}
\end{frame}


\begin{frame}{References}
\begin{reflist}
  \bibitem{g00r01} \docref{workspace-and-workflows}{The workspace and the workflow}
  \bibitem{g00r02} \docref{storage-and-sync}{Storage and sync}
  \bibitem{g00r03} \docref{gui-to-pyfs}{From the GUI to pyfs}
  \bibitem{g00r04} \docref{post-processing-definitions}{Post-processing definitions}
  \bibitem{g00r05} \recordsref
  \bibitem{g00r06} Altair Engineering, \emph{FlightStream User Manual}, version 26.1
        (build 26121), 2026. Cited with its page, as [6, p.~164].
  \bibitem{g00r07} V. Ahuja, \emph{Aerodynamic Loads over Arbitrary Bodies by Method
        of Integrated Circulation}, PhD dissertation, Auburn University, Auburn,
        2013.
  \bibitem{g00r08} D. J. Pate, B. J. German, ``A surface vorticity panel method'',
        \emph{AIAA Journal}, 56(12), 2018. DOI \texttt{10.2514/1.J057120}.
\end{reflist}
\end{frame}

\begin{frame}{References (continued)}
\begin{reflist}
  \setcounter{enumi}{8}
  \bibitem{g00r09} J. L. Hess, A. M. O. Smith, ``Calculation of potential flow about
        arbitrary bodies'', \emph{Progress in Aerospace Sciences}, 8,
        pp. 1-138, 1967.
  \bibitem{g00r10} J. Katz, A. Plotkin, \emph{Low-Speed Aerodynamics}, 2nd ed.,
        Cambridge University Press, Cambridge, 2001.
  \bibitem{g00r11} W. M. Kutta, ``Auftriebskr\"afte in str\"omenden Fl\"ussigkeiten'',
        \emph{Illustrierte Aeronautische Mitteilungen}, 6, pp. 133-135, 1902.
  \bibitem{g00r12} H. Glauert, ``The effect of compressibility on the lift of an
        aerofoil'', \emph{Proceedings of the Royal Society of London, Series
        A}, 118(779), pp. 113-119, 1928.
  \bibitem{g00r13} B. M. Simmons, S. C. Geuther, V. Ahuja, ``Validation of a
        mid-fidelity approach for aircraft stability and control
        characterization'', AIAA AVIATION Forum, 2023. DOI
        \texttt{10.2514/6.2023-4076}.
  \bibitem{g00r14} A. Gothow, F. Gavrila, J. Weiss, V. Ahuja, C. Gates,
        ``Comparative analysis: validating FlightStream with experimental
        results of a high-lift configured wing model equipped with distributed
        electric propulsion'', AIAA AVIATION Forum, 2024.
  \bibitem{g00r15} T. Theodorsen, ``General theory of aerodynamic instability and the
        mechanism of flutter'', NACA Report 496, 1935.
  \bibitem{g00r16} J. G. Leishman, \emph{Principles of Helicopter Aerodynamics}, 2nd
        ed., Cambridge University Press, Cambridge, 2006.
  \bibitem{g00r17} B. W. McCormick, \emph{Aerodynamics, Aeronautics, and Flight
        Mechanics}, 2nd ed., Wiley, New York, 1995.
\end{reflist}
\end{frame}

\begin{frame}{References (continued)}
\begin{reflist}
  \setcounter{enumi}{17}
  \bibitem{SMI} \smiref
  \bibitem{g00r19} T. Sinnige, B. Della Corte, ``TUD-XPROP: propeller geometry'',
        Delft University of Technology, Zenodo, 2026. DOI
        \texttt{10.5281/zenodo.18598046}.
  \bibitem{g00r20} \pkgauthor, ``The quasi-steady rotor against the unsteady rotor,
        measured'', report RPT-089, 2026, \code{reports/} of \url{\pkgrepo}.
  \bibitem{g00r21} \docref{getting-started}{Getting started}
  \bibitem{g00r22} \docref{fsi-workspace}{FSI in a workspace}
\end{reflist}
\end{frame}

\end{document}
